%% kapa: exam=1 round=23 year=2012 period=1 form=A subject=civil_law subject_ko=민법 seq=29 number=29
%% source: Q-Net 감정평가사 기출문제 (cst00309) artlSeq=5052578 file=raw/1008/5052578/2012년 제23회 감정평가사 제1차시험 기출문제(1교시 A형).pdf
%% conversion: kordoc 4.24.51 -> tools/gen1.py (auto); review=proofread-sonnet; flags=none
\begin{question}{29}
점유자와 회복자의 관계에 관한 설명으로 옳지 않은 것은? (다툼이 있으면 판례에 의함)
\begin{choices}
  \item 과실을 수취한 자가 선의의 점유자로 보호되기 위해서는 권원이 있다고 오신할만한 정당한 근거가 있어야 한다.
  \item 선의의 점유자가 직접 물건을 사용함으로써 얻은 이득은 과실이 아니므로 회복자에게 반환하여야 한다.
  \item 악의의 점유자는 수취한 과실의 반환뿐만 아니라 자신의 책임있는 사유로 수취하지 못한 과실의 대가도 보상하여야 한다.
  \item 선의의 타주점유자가 자신의 책임 있는 사유로 점유물을 멸실, 훼손하였다면 그로 인한 손해의 전부를 배상하여야 한다.
  \item 점유자가 과실을 취득한 경우, 통상의 필요비는 반환을 청구할 수 없다.
\end{choices}
\end{question}
\answer{2}
